\section{Pilot Human Evaluation}
\label{sec:human-eval}
%==============================================================================

This pilot asks whether human readers can tell the decoders apart at the temperature of the main experiments. It is not part of the paper's argument: KL is not used as a proxy for readability anywhere in the text.

\subsection{Design}
A blinded, self-contained offline page (\texttt{human\_\allowbreak{}eval\_\allowbreak{}tool.html}, with an English or Russian interface) has two tasks: (1) loop-detection agreement on 18 passages (clearly looping, clearly non-looping, and borderline), and (2) blind pairwise preference on 12 pairs covering all six pairings among four decoders: \textsc{dual}, SuffixMatch, and LZ Penalty, which reach 0/150 observed loops on the primary model, and FSD-style, which does not. The tool randomizes presentation order and left/right side independently per participant and ends with a copy-to-clipboard export of the raw responses (\texttt{aggregate\_\allowbreak{}human\_\allowbreak{}eval.py} combines the exports). Task 1 reports agreement with the frozen detector's classification, not with ground truth: the detector is an operational, purely syntactic instrument with a known, nonzero false-positive rate concentrated on specific corpora (Appendix~\ref{sec:event}).

\textbf{Which implementation produced the passages.} The passages were generated from the \textsc{cosine} checkpoint at $T{=}0.10$ with the decoders as they stood when the evaluation was built. The \textsc{dual}, SuffixMatch, and LZ Penalty passages are from the same code as in the tables of the main text. The FSD-style passages are from an intermediate version of the FSD-style decoder, which used fixed normalized interpolation weights $\beta^{N-k}$ for the anti-LM; it is superseded by the version of Appendix~\ref{app:fsd}, which follows the official recursive residual-mass anti-LM and is the one reported in the tables. That intermediate version reached 2/150 observed loops on the primary model, the final one 3/150. Matching the displayed text against the saved samples (\texttt{verify\_\allowbreak{}human\_\allowbreak{}eval\_\allowbreak{}passages.py}) places all six FSD-style passages in the intermediate run and none in the final run, and all passages of the other three decoders in runs that use the final code for those decoders. The FSD-style comparisons below therefore concern the intermediate version.

Nine participants were recruited informally, not through a formal panel, and sessions took a median of about 9.6 minutes (range 2--11 minutes). Each raw export was checked for internal consistency and independence before inclusion, and only sessions that passed are reported.

\begin{table}[ht]
\centering
\caption{Human evaluation, 9 participants. Task 1: loop-detection agreement against the frozen event's own classification (an operational instrument, not ground truth). Task 2: blind pairwise preference, all six pairings among the four decoders compared (\textsc{dual}, SuffixMatch, LZ Penalty, and FSD-style; the last does not reach zero loops, the other three reach 0/150); ``first'' is the method named first in each row.}
\label{tab:humaneval}
\begin{tabular}{@{}lrl@{}}
\toprule
Task 1 (loop detection) & & \\
\midrule
Pooled accuracy & 140/162 & 86.4\%, 95\% CI [78.4\%, 93.2\%] \\
\quad Looping items & 54/54 & 100.0\% \\
\quad Non-looping items & 54/54 & 100.0\% \\
\quad Borderline items & 32/54 & 59.3\% \\
Inter-participant agreement & 540/648 & 83.3\% \\
\midrule
Task 2 (pairwise preference, by participant) & & \\
\midrule
\textsc{dual} vs.\ FSD-style & 2/7 & \textsc{dual} 29\%, 95\% CI [4\%, 71\%] \\
\textsc{dual} vs.\ LZ Penalty & 3/5 & \textsc{dual} 60\%, 95\% CI [15\%, 95\%] \\
\textsc{dual} vs.\ SuffixMatch & 3/7 & \textsc{dual} 43\%, 95\% CI [10\%, 82\%] \\
FSD-style vs.\ LZ Penalty & 2/3 & FSD-style 67\%, 95\% CI [9\%, 99\%] \\
FSD-style vs.\ SuffixMatch & 2/4 & FSD-style 50\%, 95\% CI [7\%, 93\%] \\
LZ Penalty vs.\ SuffixMatch & 1/3 & LZ Penalty 33\%, 95\% CI [1\%, 91\%] \\
\bottomrule
\end{tabular}

\vspace{4pt}
{\footnotesize Task 1's CI is a participant-clustered bootstrap (10,000 resamples of the 9 participants, not the 162 individual judgments): pooling repeated judgments from the same rater as if they were independent trials (plain Clopper-Pearson) understates uncertainty, since a systematically stricter or looser rater pulls all of their own judgments the same direction. Task 2 is reported per participant (each participant's majority preference across their 2 votes per pairing collapsed to one data point, participants who split their votes evenly or tied on both excluded, hence $n\le9$ per row) with an exact (Clopper-Pearson) binomial CI, which remains valid at small $n$ and when all participants agree. See \texttt{aggregate\_\allowbreak{}human\_\allowbreak{}eval.py}'s \texttt{cluster\_\allowbreak{}level\_\allowbreak{}clopper\_\allowbreak{}pearson} (Task 2) and \texttt{participant\_\allowbreak{}cluster\_\allowbreak{}bootstrap\_\allowbreak{}ci} (Task 1).}
\end{table}

\subsection{Results}
\textbf{Loop detection} agrees with the detector's classification on 100.0\% of the unambiguously looping and of the unambiguously non-looping items and on 59.3\% of the borderline items, which were included deliberately because they probe the syntactic nature of the detector. All 22 pooled disagreements are false positives (participants calling borderline non-loops ``looping'') and none are false negatives, consistent with readers being more willing than the syntactic detector to call something a loop.

\textbf{No decoder is clearly preferred.} Across the 108 pairwise judgments (27 ties), pooled win rates are 43\% for FSD-style and SuffixMatch, 35\% for \textsc{dual}, and 30\% for LZ Penalty, and every participant-level 95\% interval in Table~\ref{tab:humaneval} includes 50\%. Individual participants differ strongly: one preferred \textsc{dual} in 6 of 12 comparisons, another called 10 of 12 pairs a tie. The data neither support nor contradict a readability cost for \textsc{dual}, and they should not be read as showing that the four decoders are equally readable. The automatic proxies (mcNLL, spelling-error rate) are not validated predictors of these judgments. This is a pilot at a single temperature on the primary model with nine participants.

\subsection{Illustrative samples}
Single samples for one prompt, at $T{=}0.10$ on the primary model, are shown only to give an impression of the text; one sample per decoder is not evidence about readability.
\begin{quote}\small
\textbf{Greedy} (loops): \emph{``they were all the subject of the consideration of the subject of the subject of the subject of the subject of the subject of the subject of the subject of the subject of the consideration of the subject of the subject of''}

\textbf{\textsc{rr-dual} $\varepsilon{=}0.05$} (no loop detected): \emph{``they had been a sensible of tranquillity and misery to take at time in terms on their sistence was always assured, after two suppose to me, with a considerably, tlength t\^ete hear or tyle so are,--''}

\textbf{LZ Penalty} (no loop detected): \emph{``they were all the conceiving which I had been so much as the stranger, and the proposal of the contrary of the present which I was always to be done the most and suffering the strong of the conversation which I had been property to the most present and all the''}
\end{quote}


%==============================================================================
